% Introduction to Part II. Unnumbered text: no \label except the \phantomsection target.
\phantomsection\label{intro:II}
Part~II extends the operators of Part~I to a background $U=e^{\ii\eta A}$ that is close to $1$ in a single global gauge: $|A|,|\partial A|,[\partial A]_\alpha\le\eps$ (hypothesis~(B), Section~\ref{sec:background}). The fields take values in the complexified Lie algebra, and the operators involve parallel transport. The constants depend in addition on $\dim\mathfrak g$ and on the decay rate $\gamma_W$ of a small potential term $W$, but not on $L$, $k$, $n$ or $U$. Part~III applies these results in a local gauge on each cube, at backgrounds that are only locally regular (Theorem~\ref{thm:G}(i)). Through Part~III they give the minimizer map and the fluctuation operator of the variational problem; Section~\ref{subsec:checklist} lists that input item by item.

\paragraph{What this Part proves.} The main result is Theorem~\ref{thm:F}. For $\eps$ small, the scalar operators $G'(U)$, $\Gamma'(U)$, $H'(U)$, $C'(U)$ and $R(U)$ satisfy the decay and pointwise bounds stated in Theorem~\ref{thm:F}(i) (Theorem~\ref{thm:cov-scalar}, Proposition~\ref{prop:cov-R}). The gauge-field Green function $\bG(U)$, built with Ba\l aban's gauge-invariant extension through the slice projection $\Pi(U)$ (Lemma~\ref{lem:slice}), satisfies the bounds of Theorem~\ref{thm:D} (Theorem~\ref{thm:cov-gauge}). The minimizer map $\bH(U)$ and the constrained Green function $\bC(U)$ satisfy the decay and pointwise bounds of Theorem~\ref{thm:E}, and have variational characterizations (Theorem~\ref{thm:cov-min}). All these operators are analytic in $A$ and $W$ and covariant under $G$-valued gauge transformations, and the identities $\bG DR=DG'R$, $RD^T\bG=RG'D^T$ and $\bQ_\pi\bG DR=0$ are exact. Proposition~\ref{prop:covariance} transfers the bounds to every background obtained from one that satisfies~(B) by a $G$-valued gauge transformation, with covariant derivatives, the H\"older bounds in the gauge of~(B); the comparisons with the operators at $U=1$ hold only in the gauge of~(B). The proofs write every operator as its counterpart at $U=1$ plus an admissible perturbation (Definition~\ref{def:admissible}). Combes--Thomas, inversion with decay and local regularity are stable under such perturbations (Lemmas~\ref{lem:CT-stable}, \ref{lem:perturbed-inverse} and~\ref{lem:bootstrap}).

\paragraph{Which averaging.} Part~II is stated for a model averaging pair $(Q'(U),\bQ(U))$, with transport along one coordinate-ordered contour and only the main term of Ba\l aban's linearized averaging. It intertwines the covariant derivatives up to a defect: $\bQ(U)D_U=D_VQ'(U)+E(U)$, with $\|E(U)\|\le C\eps$ (Lemma~\ref{lem:cov-avg}(c)). The model pair is a proof device. For the linearization of Ba\l aban's symmetric averaging $E=0$, and Corollary~\ref{cor:parts-apply} applies the real results of Parts~II and~III to it, through Remark~\ref{rem:general-averaging}. A first reading may set $E=0$.

\paragraph{Imports.} From Part~I: the operators at $U=1$ (Theorems~\ref{thm:scalar-green} and~\ref{thm:gauge-green}, Propositions~\ref{prop:C-scalar}, \ref{prop:R} and~\ref{prop:QGQ-gauge}), Lemmas~\ref{lem:avg} and~\ref{lem:decay-calculus}--\ref{lem:CT-form}, and Theorem~\ref{thm:local-regularity}. No proof in Part~II uses a later result.

\paragraph{Exports.} Part~III uses all of Part~II. The expansion of covariant derivatives and the curvature operator (Lemmas~\ref{lem:cov-expansion} and~\ref{lem:curvature}), the perturbed inverses and the bootstrap (Lemmas~\ref{lem:perturbed-inverse} and~\ref{lem:bootstrap}), and the slice projection (Lemma~\ref{lem:slice}) are used throughout Parts~III--X. Theorem~\ref{thm:cov-min} is used again in Parts~V--IX and~XI. All of Part~II is on the main line. Section~\ref{subsec:checklist} is a list, not a result: it maps the input of Ba\l aban's variational theorem to the places where this note proves it, mostly in Parts~III--VI.

\paragraph{Details for this Part.} The companion has no chapter on Part~II; its parameter ledger lists the constants of Part~II that Parts~VII--XI use.
